% GDZ PPN236018647,native592–594through129;HKUST593–595.
\clearpage\SourceLeaf{592}\ArithmeticHead{563}{R}{1801 SEPT. 7. MEDIO—1805 AUG. 30.}
\NumberedItem{121}
\noindent\hspace*{\parindent}Formulas permultas novas in Astronomia Theorica utilissimas eruimus.
\par\noindent\hfill 1801 Mense Octobr.
\par\medskip\begingroup\small
\noindent\hspace*{\parindent}Siehe die beiden vorhergehenden Eintragungen, sowie die folgende.
\par\noindent\hfill\textsc{Brendel.}
\par\endgroup

\NumberedItem{122}
\noindent\hspace*{\parindent}Annis insequentibus 1802, 1803, 1804 occupationes astronomicae maximam otii partem abstulerunt, calculi imprimis circa planetarum novorum theoriam instituti. Unde evenit, quod hisce annis catalogus hicce neglectus est. Dies itaque, quibus aliquid ad matheseos incrementa conferre datum est, memoriae exciderunt.
\par\medskip\begingroup\small
\noindent\hspace*{\parindent}Im November 1801 hatte \textsc{Gauss} mehrere Bahnen der Ceres berechnet, die \textsc{Piazzis} Beobachtungen (Vergl. die Bemerkung zu der Nr. 120) hinreichend gut darstellten, und die nebst einer Ephemeride zur Aufsuchung unter der Überschrift \textit{Fortgesetzte Nachrichten über den längst vermutheten neuen Haupt-Planeten unseres Sonnen-Systems} in der Monatlichen Correspondenz zur Beförderung der Erd- und Himmelskunde, Dezember 1801 (Werke VI, S. 199), veröffentlicht wurden. Die Rechnungen nebst den Rechnungsvorschriften finden sich ganz oder doch zum größten Teil in den bei der Nr. 119 genannten Schedae Ag und Ah. Hiernach wurde Ceres auch am 7. Dezember von \textsc{von Zach} aufgefunden und ihre Identität durch eine weitere Beobachtung von \textsc{Olbers} am 1. Januar 1802 bestätigt.
\par\noindent\hspace*{\parindent}\textsc{Gauss} benutzte die neuen Beobachtungen sogleich zu einer genaueren Bahnbestimmung und vertiefte sich auf lange Zeit, auch noch über das Jahr 1804 hinaus, so in die Bahnberechnungen der Ceres und der später entdeckten Pallas, daß er seine übrigen Untersuchungen fast ganz liegen ließ. Die Schedae Ai, Ak, Al sind voll von solchen Berechnungen. Im Jahre 1802 beschäftigte er sich außerdem mit der Berechnung der Störungen der Ceres, welche Werke VII, 1906, S. 377 ff. abgedruckt sind.
\par\noindent\hspace*{\parindent}An der eigentlichen Verbesserung seiner Methode zur Bahnbestimmung scheint \textsc{Gauss} auch in dieser Zeit weniger gearbeitet zu haben, denn aus den Eintragungen Nr. 125 bis Nr. 127 ist zu schließen, daß die Form, in der diese Methode uns in der \textit{Theoria motus} entgegentritt, erst im Jahre 1806 entstanden ist.
\par\noindent\hspace*{\parindent}Man vergleiche \textsc{Gauss}' Briefwechsel mit \textsc{Olbers} und Werke VI.
\par\noindent\hfill\textsc{Brendel.}
\par\endgroup

\NumberedItem{123}
\begingroup\renewcommand{\thefootnote}{[*)]}
\noindent\hspace*{\parindent}Demonstratio theorematis venustissimi supra 1801 Mai.\footnotemark\ commemorati, quam per 4 annos et ultra omni contentione quaesiveramus, tandem perfecimus. Comment[ationes] rec[entiores], I.
\renewcommand{\thefootnote}{[*)}
\footnotetext{Nr. 118, oben S. 560.]}
\par\noindent\hfill 1805 Aug. 30.
\par\endgroup
\par\medskip\begingroup\small
\noindent\hspace*{\parindent}Es handelt sich um die Untersuchungen der im I. Bande der Commentationes societatis regiae \mbox{scien-}
\par\endgroup
\par\medskip\noindent\hfill 71\textsuperscript{*}
\clearpage\SourceLeaf{593}\ArithmeticHead{564}{L}{TAGEBUCH MIT ERLÄUTERUNGEN.}
\par\begingroup\small
\noindent tiarum Gottingensis recentiores veröffentlichten Abhandlung \textit{Summatio quarumdam serierum singularium} (vorgelegt am 24. August 1808), Werke II, S. 9, siehe auch die Anzeige ebenda S. 155. Eine lebhafte Schilderung seiner fortgesetzten, lange vergeblich gebliebenen Bemühungen um den Beweis und seines endlichen Erfolgs gibt \textsc{Gauss} in dem oben S. 24 abgedruckten Briefe an \textsc{Olbers} vom 3. September 1805. Vergl. den Artikel 18 des \textsc{Bachmann}schen Aufsatzes, Werke X2, S. 45.
\par\noindent\hfill\textsc{Klein.}
\par\endgroup

\NumberedItem{124}
\noindent\hspace*{\parindent}Theoriam interpolationis ulterius excoluimus.
\par\noindent\hfill 1805 Novbr.
\par\medskip\begingroup\small
\noindent\hspace*{\parindent}Das im Oktober 1805 begonnene Handbuch, betitelt \textit{Mathematische Brouillons} (18, Bd), wird mit einer Aufzeichnung über Interpolation eröffnet, die (vergl. die Bemerkung von \textsc{Schering} Werke III, S. 328) als ein erster Entwurf der Werke III, S. 265 aus dem Nachlaß abgedruckten Abhandlung \textit{Theoria interpolationis methodo nova tractata} anzusehen ist. Daraus folgt im Einklang mit unserer Tagebuchnotiz, daß diese Abhandlung nicht vor dem November 1805 verfaßt sein kann. Andererseits ist sie aber jedenfalls vor dem 25. August 1806 geschrieben, weil bei dem im art. 41. (Werke III, S. 325) gegebenen Beispiel die Exzentrizität der Juno entsprechend den V. Elementen dieses kleinen Planeten gleich \(0{,}254236\) genommen wird, vergl. die auf eben diesen Wert bezügliche Bemerkung, oben S. 443.
\par\noindent\hfill\textsc{Klein.\quad Schlesinger.}
\par\endgroup

\NumberedItem{125}
\noindent\hspace*{\parindent}Methodum ex duobus locis heliocentricis corporis circa solem moventis eiusdem elementa determinandi novam perfectissimam deteximus.
\par\noindent\hfill 1806 Ianuar.
\par\medskip\begingroup\small
\renewcommand{\thefootnote}{*)}
\noindent\hspace*{\parindent}\textsc{Gauss} berichtet hierüber in dem Briefe an \textsc{Olbers} vom 3. Februar 1806 (\textit{Wilhelm Olbers, Sein Leben und seine Werke} II1, Berlin 1900, S. 287); man vergleiche die Bemerkungen zu den Nummern 120 und 122 und die artt. 88.—97. der \textit{Theoria motus} (Werke VII, 1906, S. 112), sowie den Bericht\footnote{\textit{Über die Aufgabe: »Aus zwey ihrer Grösse und Lage nach gegebenen Radii Vectores und der verflossenen Zeit die elliptischen Elemente der Planetenbahn zu bestimmen«}. Nach § 88—97 der \textit{Theoria motus corporum coelestium etc. etc.} des Hrn. Prof. \textsc{Gauss}. [Der ungenannte Verfasser ist vermutlich \textsc{v. Zach}.]} über die letzteren in der Monatlichen Correspondenz \textbf{20}, 1809, S. 322. — Im Nachlaß finden sich die entsprechenden Untersuchungen im Handbuch 18, Bd, S. 52 und an anderen Stellen. Bei seinen früheren Bahnbestimmungen hat sich \textsc{Gauss} einer weniger vollkommenen Methode bedient.
\par\noindent\hfill\textsc{Brendel.}
\par\endgroup

\NumberedItem{126}
\noindent\hspace*{\parindent}Methodum e tribus planetae locis geocentricis eius orbitam determinandi ad summum perfectionis gradum eveximus.
\par\noindent\hfill 1806 Mai.
\clearpage\SourceLeaf{594}\ArithmeticHead{565}{R}{1805 AUG. 30—1807 FEBR. 17.}
\par\begingroup\small
\noindent\hspace*{\parindent}Man vergleiche \textit{Theoria motus}, liber secundus, artt. 115—163. (Werke VII, 1906, S. 155). Im Nachlaß stehen die entsprechenden Untersuchungen im Handbuch 18, Bd, S. 94.
\par\noindent\hfill\textsc{Brendel.}
\par\endgroup

\NumberedItem{127}
\noindent\hspace*{\parindent}Methodus nova ellipsin et hyperbolam ad parabolam reducendi.
\par\noindent\hfill 1806 April.
\par\medskip\begingroup\small
\noindent\hspace*{\parindent}Man vergleiche \textit{Theoria motus}, artt. 33 ff., Werke VII, 1906, S. 45, sowie im Nachlaß das Handbuch 18, Bd, S. 70.
\par\noindent\hfill\textsc{Brendel.}
\par\endgroup

\NumberedItem{128}
\noindent\hspace*{\parindent}Eodem circiter tempore resolutionem functionis \(\frac{x^p-1}{x-1}\) in factores quatuor absolvimus.
\par\noindent\hfill [1806 April.—Mai.]
\par\medskip\begingroup\small
\noindent\hspace*{\parindent}Im art. 22. der am 5. April 1825 der Gesellschaft der Wissenschaften vorgelegten \textit{Theoria Residuorum biquadraticorum Commentatio prima}, Werke II, S. 89 sagt \textsc{Gauss} in bezug auf die hier angezeigte Zerlegung, daß sie im engsten Zusammenhang stehe mit den in den artt. 15.—20. jener Abhandlung enthaltenen Untersuchungen, und daß sie sich mit Hilfe dieser Untersuchungen ohne Schwierigkeit vollständig erledigen lasse. Sed — fährt er fort — hanc tractationem ad aliam occasionem nobis reservamus. Im Nachlaß hat sich keine auf diesen Gegenstand bezügliche Aufzeichnung vorgefunden.
\par\noindent\hfill\textsc{Bachmann.}
\par\endgroup

\NumberedItem{129}
\noindent\hspace*{\parindent}Methodus nova e quatuor planetae locis geocentricis, quorum duo extremi sunt incompleti, eius orbitam determinandi.
\par\noindent\hfill 1807 Ian. 21.
\par\medskip\begingroup\small
\noindent\hspace*{\parindent}Man vergleiche den Brief von \textsc{Gauss} an \textsc{Olbers} vom 27. Januar 1807, \textit{Wilhelm Olbers, Sein Leben und seine Werke} II1, Berlin 1900, S. 320, sowie \textit{Theoria motus}, liber secundus, artt. 164.—171., Werke VII, 1906, S. 222, ferner im Nachlaß das Handbuch 18, Bd, S. 133.
\par\noindent\hfill\textsc{Brendel.}
\par\endgroup
